\documentclass[12pt,a4paper]{article}

\usepackage[margin=1in]{geometry}
\usepackage{graphicx}
\usepackage{float}
\usepackage{booktabs}
\usepackage{array}
\usepackage{hyperref}
\usepackage{microtype}
\usepackage[T1]{fontenc}
\usepackage[utf8]{inputenc}

\hypersetup{colorlinks=true,linkcolor=black,urlcolor=blue}
\setlength{\parskip}{0.45em}
\setlength{\parindent}{0pt}
\renewcommand{\arraystretch}{1.15}

\newcommand{\evidenceimage}[1]{%
  \includegraphics[width=0.96\linewidth,height=0.72\textheight,keepaspectratio]{\detokenize{#1}}%
}

\begin{document}

% Cover page supplied for this project. Add the official logo file at the
% project root to replace the text fallback below.
\thispagestyle{empty}
\begin{center}
\vspace*{0.5cm}
{\Large \textbf{Tank-Full Smart Dehumidifier Controller}}\\
\vspace{0.8cm}
{\large \textbf{Digital System Design (EE311)}}\\
\vspace{1.2cm}
A PROJECT REPORT SUBMITTED IN PARTIAL\\
FULFILLMENT OF THE REQUIREMENTS FOR THE AWARD\\
OF THE DEGREE
\vspace{0.6cm}

{\large \textbf{BACHELOR OF TECHNOLOGY}}\\
\vspace{0.2cm}
IN\\
\vspace{0.2cm}
{\large \textbf{Electrical and Electronics Engineering}}\\
\vspace{1.2cm}
Submitted by:\\
\vspace{0.3cm}
\textbf{Adarsh Akshat (231EE103)}\\
\vspace{0.8cm}
Under the supervision of\\
\vspace{0.2cm}
\textbf{Dr. Yashwant Kashyap}\\
\vspace{0.8cm}
\IfFileExists{nitk_logo.jpg}{%
  \includegraphics[width=0.28\textwidth]{nitk_logo.jpg}%
}{%
  \fbox{\parbox[c][1.5cm][c]{0.28\textwidth}{\centering NITK logo}}%
}
\vfill
{\large \textbf{DEPARTMENT OF ELECTRICAL AND ELECTRONICS ENGINEERING}}\\
\vspace{0.2cm}
\textbf{NATIONAL INSTITUTE OF TECHNOLOGY KARNATAKA, SURATHKAL}\\
\vspace{0.1cm}
SRINIVASNAGAR, MANGALORE -- 575025\\
\vspace{0.3cm}
\textbf{November 2026}
\end{center}

\clearpage
\pagenumbering{roman}
\tableofcontents
\clearpage
\pagenumbering{arabic}

\section{Abstract}

This report describes a synchronous digital controller for a tank-full smart dehumidifier demonstration on the Digilent Basys 3 FPGA board. The design models humidity demand with an adjustable percentage value, applies hysteresis and a minimum compressor off-time, gives tank-full and fault conditions priority, and provides defrost and post-defrost recovery states. A multiplexed seven-segment display and LEDs expose the value and controller status. The report documents the circuit view, design methodology, finite-state machine, timing behavior, simulation evidence, FPGA implementation reports, and hardware demonstration.

The current implementation is a digital control demonstrator. It does not sample a physical humidity or tank sensor and does not drive an actual compressor or fan. Push buttons provide the demo inputs, and the modeled outputs are shown on LEDs.

\section{Objectives and Scope}

The project objectives are to:
\begin{itemize}
  \item implement modular, synchronous Verilog control logic for a dehumidifier operating sequence;
  \item demonstrate button synchronization and debouncing, humidity hysteresis, tank-full shutdown, defrost/recovery, and fault latching;
  \item display the demo humidity value and operating status on the Basys 3 board; and
  \item provide simulation and FPGA implementation evidence for the design.
\end{itemize}

Humidity is a user-adjustable demo value, changed in increments of five percentage points. Tank status is advanced through Normal, Low, and Full by a button. These signals model sensor information for this project. The RTL signals named \texttt{compressor\_on} and \texttt{fan\_on} are logic-level status outputs; an electrical power stage is not part of the implementation.

\section{System Architecture}

The top-level module, \texttt{top\_controller}, connects the button debouncers, humidity and tank controllers, dehumidifier state machine, display driver, and LED mapper. The design is clocked from the Basys 3 100 MHz oscillator. The supplied XDC targets part \texttt{xc7a35tcpg236-1} and maps the clock, five push buttons, fifteen LEDs, seven segment lines, and four digit enables.

\subsection{Functional modules}
\begin{itemize}
  \item \texttt{button\_debouncer}: two-flop input synchronizer, stable-level filter, and one-cycle press event.
  \item \texttt{humidity\_controller}: resettable value adjusted in 5\% steps and limited to 0--100\%.
  \item \texttt{tank\_controller}: cycles through Normal, Low, and Full status values.
  \item \texttt{dehumidifier\_controller}: humidity hysteresis, state sequencing, safety output gates, fault latch, anti-cycle counter, and defrost recovery timer.
  \item \texttt{seven\_segment\_display}: multiplexes the three decimal digits and blanks the leftmost digit.
  \item \texttt{led\_status\_display}: maps humidity and operating flags to the fifteen board LEDs.
\end{itemize}

\subsection{Board input and output mapping}
\begin{center}
\begin{tabular}{>{\raggedright\arraybackslash}p{0.21\linewidth}p{0.69\linewidth}}
\toprule
Input / output & Function \\
\midrule
\texttt{btnC} & Reset all modeled control state \\
\texttt{btnU}, \texttt{btnD} & Increase or decrease demo humidity by 5\% \\
\texttt{btnL} & Cycle tank status: Normal $\rightarrow$ Low $\rightarrow$ Full $\rightarrow$ Normal \\
\texttt{btnR} & Toggle defrost request \\
\texttt{led[6:0]} & Humidity value in binary \\
\texttt{led[7]}, \texttt{led[8]} & Modeled compressor and fan outputs \\
\texttt{led[9:14]} & Dehumidify, defrost, tank-full, fault, tank-low, and anti-cycle flags \\
\texttt{seg}, \texttt{an} & Active-low seven-segment patterns and digit enables \\
\bottomrule
\end{tabular}
\end{center}

\section{Design Methodology}

After initialization, button inputs are synchronized and debounced. A valid up or down event changes the humidity demo value by five percentage points, while the tank button advances its three-state sequence and the defrost button toggles the defrost request. On every clock cycle, the control logic updates the humidity demand using hysteresis and evaluates safety and operating conditions. The display and LEDs continuously reflect the resulting value and state.

The controller gives an invalid tank encoding highest priority and latches a fault. A full tank gates both modeled compressor and fan outputs off. Defrost turns the compressor off while retaining fan operation. On leaving defrost, the controller enters a timed recovery interval before it may return to dehumidification. Normal humidity demand is accepted only when the minimum off-time has expired.

\begin{figure}[H]
\centering
\evidenceimage{Tank-Full Dehumidifier Control Flowchart.png}
\caption{Methodology flowchart. Sensor-related labels in this conceptual diagram correspond to button-controlled demo inputs in the implemented RTL.}
\end{figure}

\section{Finite-State Machine}

The dehumidifier controller uses six states. The encoded values are defined in \texttt{dehumidifier\_controller.v}.

\begin{center}
\begin{tabular}{lll}
\toprule
State & Code & Function \\
\midrule
IDLE & 000 & Wait for a humidity demand and expired off-time \\
DEHUMIDIFY & 001 & Modeled compressor and fan enabled \\
DEFROST & 010 & Compressor disabled; fan may operate \\
RECOVERY & 011 & Timed post-defrost wait \\
TANK\_FULL & 100 & Stop modeled compressor and fan while tank is full \\
FAULT & 101 & Latched safe state for invalid tank status \\
\bottomrule
\end{tabular}
\end{center}

Demand sets at or above 60\% and clears at or below 50\%; between the thresholds the prior demand is retained. The implementation defaults to an 18,000,000,000-cycle minimum off-time and a 500,000,000-cycle recovery interval. At 100 MHz these correspond to 180 seconds and 5 seconds, respectively. A reset clears the fault latch. Safety conditions gate the modeled outputs immediately, in addition to advancing the registered state machine.

\begin{figure}[H]
\centering
\evidenceimage{Dehumidifier Controller FSM Diagram.png}
\caption{Finite-state machine diagram and state encoding.}
\end{figure}

\section{Circuit Diagram and FPGA Mapping}

The following Vivado schematic shows the synthesized top-level controller hierarchy and its connections to the FPGA I/O. It serves as the circuit-level view of the RTL implementation. The design uses logic inside the Artix-7 FPGA; the external appliance sensors and power drivers are not included.

\begin{figure}[H]
\centering
\evidenceimage{Schematic.png}
\caption{Synthesized RTL schematic for the dehumidifier controller.}
\end{figure}

The Basys 3 constraints specify a 10 ns clock period. The default display refresh interval is 25,000 clock cycles (250 microseconds per digit), giving an approximately 1 kHz complete four-digit scan. The default button filter interval is 1,000,000 cycles (10 ms).

\section{Timing Behavior}

The timing diagram illustrates reset, a debounced humidity-up event, the 55\% to 60\% change, activation of dehumidification, and shutdown when tank-full status is asserted. The illustration is qualitative and is not drawn to scale. Exact state and signal transitions are determined by the RTL and its clocked counters.

\begin{figure}[H]
\centering
\evidenceimage{Basys 3 Dehumidifier Timing Diagram.png}
\caption{Illustrative timing diagram of the input, control, and status behavior.}
\end{figure}

\section{Simulation and Verification}

The self-checking testbench is \texttt{tb/tb\_top\_controller.v}. It instantiates the design with shortened debounce, lockout, recovery, and display-refresh parameters so that the scenarios complete quickly in simulation. The checks cover startup and reset, short and long button presses, humidity changes, seven-segment digits, hysteresis, anti-cycle lockout, tank-low/full behavior, defrost and recovery, invalid tank status and fault reset, and LED mapping. The testbench emits a PASS/FAIL summary when the run completes.

The supplied waveform screenshot shows the early simulation interval, with five checks completed and no failures at the displayed cursor. It is a partial waveform capture; the testbench's final transcript is the source for the complete run summary.

\begin{figure}[H]
\centering
\evidenceimage{Simulation_waveform.png}
\caption{Simulation waveform capture, including button inputs, LED/display outputs, and testbench counters.}
\end{figure}

\section{FPGA Implementation Results}

The supplied Vivado utilization capture reports 162 slice LUTs and 220 slice registers for the top-level design. It also reports 32 bonded I/O blocks and one global clock buffer. The timing summary capture reports a worst negative slack of 4.368 ns, total negative slack of 0.000 ns, and zero failing setup endpoints. Hold slack is 0.139 ns with zero failing hold endpoints. These values describe the captured implementation run and should be regenerated if the source, constraints, or Vivado version changes.

\begin{figure}[H]
\centering
\evidenceimage{Utilisation_Report.png}
\caption{Vivado utilization report for the captured design run.}
\end{figure}

\begin{figure}[H]
\centering
\evidenceimage{Timing_Reports.png}
\caption{Vivado design timing summary for the captured implementation.}
\end{figure}

The project overview and implemented-design views provide further evidence of the selected top module, constraints, and Basys 3 target.

\begin{figure}[H]
\centering
\evidenceimage{Project_overview.png}
\caption{Vivado project hierarchy showing the RTL, constraints, and simulation testbench.}
\end{figure}

\begin{figure}[H]
\centering
\evidenceimage{Implemented_design.png}
\caption{Vivado implemented-design view for the selected Artix-7 device.}
\end{figure}

\section{Hardware Demonstration}

The project was demonstrated on a Basys 3 board. The photograph shows the board displaying a humidity value and status LEDs. This validates the board-level demonstration of the FPGA logic and display outputs; it does not represent an electrical connection to a working dehumidifier.

\begin{figure}[H]
\centering
\evidenceimage{Working_pic.jpeg}
\caption{Basys 3 board during the controller demonstration.}
\end{figure}

\begin{figure}[H]
\centering
\evidenceimage{basys3_hardware_walkaround.png}
\caption{Basys 3 board features used to identify the buttons, display, and LEDs.}
\end{figure}

\section{Limitations and Future Scope}

The current build uses push buttons to model humidity and tank state. It does not contain sensor interfaces, environmental calibration, a compressor or fan driver, or a mains safety circuit. The LED outputs indicate modeled logic states only. The timing and utilization captures apply to the current source and constraints.

Future work includes:
\begin{itemize}
  \item adding an ILA core and capturing FSM, demand, tank, timer, and output signals for live debug; the setup wizard screenshot is \texttt{Debug\_ILA(further).png};
  \item exploring appropriate Vivado IP catalog components and recording any integration; the design-view reference is \texttt{IP\_Catalog(further).png};
  \item integrating calibrated humidity and tank-level sensors with suitable input conditioning; and
  \item designing and validating isolated, appropriately rated drivers before connecting any real compressor or fan.
\end{itemize}

The two files marked \texttt{(further)} are included as future-work references and are not represented as features already implemented in the RTL.

\begin{figure}[H]
\centering
\evidenceimage{Debug_ILA(further).png}
\caption{Vivado ILA setup wizard retained as a reference for future on-chip debug instrumentation. No ILA core is instantiated in the current RTL.}
\end{figure}

\begin{figure}[H]
\centering
\evidenceimage{IP_Catalog(further).png}
\caption{Vivado implemented-design view retained for further design and IP exploration.}
\end{figure}

\section{Conclusion}

The project implements a modular FPGA control demonstrator with user-adjustable humidity, hysteresis, anti-cycle timing, tank-full shutdown, defrost and recovery, fault handling, and board-level visual status. The included schematic, methodology flowchart, FSM, timing diagram, simulation waveform, implementation reports, and hardware photographs document the design and its demonstration. Sensor acquisition and appliance power control remain future extensions.

\end{document}
